\chapter{Conclusiones y Recomendaciones}

\section{Conclusiones}
A partir del desarrollo e implementación del Asistente Virtual Inteligente, se han establecido las siguientes conclusiones respecto al cumplimiento de los objetivos planteados:

\begin{itemize}
    \item El análisis de los procesos actuales permitió identificar limitaciones asociadas con la consulta manual de documentos y la ausencia de mecanismos de búsqueda semántica, lo cual permitió establecer los requerimientos funcionales y no funcionales del prototipo.

    \item Se diseñó una arquitectura modular cliente-servidor que integró el motor RAG, los servicios de voz, el módulo de interacción visual y el avatar 3D. El uso de WebSockets permitió implementar comunicación bidireccional y procesamiento asíncrono entre cliente y backend.

    \item Se implementó el motor RAG utilizando el modelo de embeddings \texttt{paraphrase-multilingual-MiniLM-L12-v2}, obteniendo un \textit{Hit Rate} Top-3 global de 92,0\% y un MRR de 0,84 en el conjunto de consultas empleado para la evaluación. Estos resultados evidencian un desempeño favorable del mecanismo de recuperación dentro de las condiciones evaluadas.

    \item La integración del avatar 3D, reconocimiento de voz, síntesis de voz y mecanismos de visión artificial permitió obtener un prototipo multimodal funcional. Las mediciones realizadas evidenciaron tiempos compatibles con una interacción conversacional en el entorno de prueba, aunque el rendimiento final depende de la infraestructura y de los servicios externos empleados.

    \item La evaluación piloto realizada con 13 estudiantes presentó una tendencia favorable en la percepción de integración, utilidad e intención futura de uso del prototipo. Debido al carácter exploratorio y al tamaño de la muestra, estos resultados se interpretan como evidencia inicial de aceptación y no como resultados generalizables a toda la población estudiantil.
\end{itemize}

\section{Recomendaciones}
Con base en los retos enfrentados y el análisis técnico del sistema en producción, se presentan las siguientes recomendaciones:

\begin{itemize}
    \item \textbf{Escalabilidad y Pruebas de Estrés Exhaustivas:} Dado que las pruebas de concurrencia mediante JMeter tuvieron un alcance exploratorio para validar el funcionamiento asíncrono, se recomienda realizar pruebas de estrés a gran escala antes de un despliegue institucional definitivo. Esto permitirá dimensionar adecuadamente la infraestructura requerida e implementar estrategias de balanceo de carga en entornos en la nube (cloud) cuando el sistema migre desde el servidor local de desarrollo.

    \item \textbf{Dependencia de Red y Latencia (Limitación de APIs de Terceros):} Dado que los tiempos de respuesta de la transcripción y síntesis de voz dependen de servicios en la nube externos (ElevenLabs y Groq), el sistema es susceptible a variaciones en el ancho de banda institucional. Se recomienda explorar o implementar localmente modelos de voz (\textit{On-Premise}) o sistemas de caché para mitigar retrasos si la conexión a internet de la universidad presenta intermitencias.
    
    \item \textbf{Calidad de la Base Documental (Limitación de Recuperación):} El motor RAG demostró mitigar efectivamente las alucinaciones siempre y cuando los documentos indexados sean de alta legibilidad semántica. Se recomienda que la administración del sistema aplique un protocolo estricto de curaduría de archivos PDF antes de cargarlos, rechazando documentos escaneados sin Reconocimiento Óptico de Caracteres (OCR) válido, ya que degradan la capacidad de búsqueda.
    
    \item \textbf{Control de costos operativos:} Debido a que los servicios externos de voz pueden incrementar su costo en función del volumen de solicitudes, se recomienda monitorear el consumo real durante una eventual implementación institucional y evaluar alternativas locales o esquemas de control de uso cuando sean necesarios.
\end{itemize}

\section{Trabajos Futuros}
El prototipo actual sienta una base tecnológica sólida sobre la cual pueden derivarse múltiples líneas de investigación o ampliación de producto:

\begin{itemize}
    \item \textbf{Integración Profunda con el SGA Institucional:} Una evolución natural del proyecto consistiría en integrar la sesión autenticada del estudiante con las bases de datos transaccionales de la universidad (mediante APIs RESTful). Esto permitiría al asistente no solo informar sobre normas generales, sino brindar respuestas hiper-personalizadas, tales como "Dada tu nota actual, te sugiero tomar estas dos materias como prerrequisito".
    
    \item \textbf{Generación de Emociones Dinámicas (Affective Computing):} Actualmente el avatar realiza sincronización labial y animaciones programadas. Como trabajo futuro, se puede entrenar un modelo de análisis de sentimientos que evalúe la carga emocional de la respuesta del LLM en tiempo real, activando \textit{blendshapes} específicos (sonrisas ante buenas noticias o expresiones de empatía ante problemas) para dotar al avatar de inteligencia emocional.
    
    \item \textbf{Análisis Longitudinal del Impacto Estudiantil:} Ejecutar el protocolo de Google Forms diseñado a lo largo de uno o dos periodos académicos completos para analizar estadísticamente si la disponibilidad continua del asistente reduce el abandono estudiantil provocado por la desinformación de procesos administrativos.
\end{itemize}
